% Numbers from results/tables/final_fairness_rolling.csv (notebook 03, Section 3.4).
\begin{itemize}
  \item \textbf{Con memoria, la desigualdad se reduce mes a mes} (Figura~\ref{fig:rolling}). Con el modelo, la diferencia entre el más y el menos cargado pasa de 21,7 al principio a 14,8 tras enero y a 3,5 tras mayo, y la desviación media baja de 6,3 a 0,5. No cuesta nada: en los cinco meses se cubren todas las guardias ($f_1 = 0$) y se respetan todas las peticiones ($f_2 = 0$).
  \item \textbf{Sin memoria no se corrige nada.} Equilibrar cada mes por separado deja la desigualdad acumulada igual o peor que no tener ningún criterio de equidad: tras mayo el rango es 26,1 sin memoria y 23,7 sin criterio, frente a 3,5 con memoria. Repartir bien los festivos de cada mes da los mismos festivos a quien ya arrastraba un exceso que a quien iba por detrás: \emph{la memoria es lo que hace justo el reparto}.
  \item \textbf{La métrica importa poco a medio plazo.} La desviación media reparte mejor entre los residentes intermedios al principio (4,0 frente a 4,4 tras enero) y el rango protege mejor al más cargado (14,8 frente a 16,0), pero tras cinco meses ambas llegan al mismo reparto (rango 3,5). Recomendamos mantener el rango: es más fácil de explicar («la diferencia entre quien más y quien menos festivos lleva») y se resuelve en menos de un segundo, mientras que la desviación media tarda hasta 36~s en algunos meses.
  \item \textbf{El rango de un solo mes engaña.} En enero el rango se queda en 14,8 aunque cambie la plantilla (Sección~\ref{sec:plantilla}), porque lo fija el residente con más histórico. Visto a lo largo de varios meses el criterio sí funciona: lo que no se puede compensar en un mes se compensa en los siguientes.
\end{itemize}

{\small Detalle técnico: para esta simulación los niveles del objetivo se resuelven uno a uno, con 20~s cada uno. La política sin memoria agota ese límite en el nivel de equidad (tiene muchas soluciones simétricas); las demás se resuelven hasta optimalidad probada.\par}
